%=============================================================================
% 总结与作业
%=============================================================================

\section{本周回顾}

% -----------------------------------------------------------------------------
% 第1页：预处理流程回顾
% -----------------------------------------------------------------------------

\begin{frame}{本周知识体系：图像预处理流水线}
    \begin{center}
    \begin{tikzpicture}[
        box/.style={draw, rounded corners=4pt, minimum width=2.2cm, minimum height=1cm,
                     align=center, font=\small\bfseries, fill=#1},
        sub/.style={font=\tiny, align=center, text=gray!80!black},
        >=stealth
    ]
        % 节点 - 用绝对坐标控制间距
        \node[box=blue!15]   (input)  at (0, 0)    {原始图像};
        \node[box=green!20]  (denoise) at (2.8, 0)  {去噪};
        \node[box=yellow!25] (enhance) at (5.6, 0)  {增强};
        \node[box=orange!20] (binary)  at (8.4, 0)  {二值化};
        \node[box=red!15]    (geo)     at (11.2, 0) {几何矫正};
        \node[box=purple!15] (output)  at (14, 0)   {预处理\\结果};

        % 箭头
        \draw[->, thick] (input)   -- (denoise);
        \draw[->, thick] (denoise) -- (enhance);
        \draw[->, thick] (enhance) -- (binary);
        \draw[->, thick] (binary)  -- (geo);
        \draw[->, thick] (geo)     -- (output);

        % 关键函数标注
        \node[sub, below=6pt of denoise] (d_sub) {GaussianBlur\\medianBlur};
        \node[sub, below=6pt of enhance] (e_sub) {CLAHE\\equalizeHist};
        \node[sub, below=6pt of binary]  (b_sub) {threshold\\Otsu\\adaptive};
        \node[sub, below=6pt of geo]     (g_sub) {warpPerspective\\warpAffine};

        % Week3 虚线框（包含函数标注）
        \node[draw=blue!60, dashed, rounded corners=6pt,
              inner sep=8pt, fit=(denoise)(enhance)(binary)(geo)
                                     (d_sub)(e_sub)(b_sub)(g_sub),
              label={[blue!70, font=\small\bfseries]above:Week 3 本周所学}] {};
    \end{tikzpicture}
    \end{center}

    \vspace{0.2cm}
    \begin{alertblock}{核心思想}
        预处理的目标：将"脏"的原始试卷图像 $\longrightarrow$ "干净"的标准图像，为后续 OCR 和版面分析做准备。
    \end{alertblock}
\end{frame}

% -----------------------------------------------------------------------------
% 第2页：核心函数回顾
% -----------------------------------------------------------------------------

\begin{frame}[fragile]{核心函数回顾：本周学到哪些API？}
    \begin{columns}
        \column{0.5\textwidth}
        \begin{block}{去噪函数}
            \begin{lstlisting}[language=Python, basicstyle=\ttfamily\tiny]
# 高斯滤波
blur = cv2.GaussianBlur(
    gray, (5, 5), 0)

# 中值滤波（去椒盐噪声）
median = cv2.medianBlur(gray, 5)
            \end{lstlisting}
        \end{block}

        \begin{block}{增强函数}
            \begin{lstlisting}[language=Python, basicstyle=\ttfamily\tiny]
# CLAHE自适应对比度增强
clahe = cv2.createCLAHE(
    clipLimit=2.0,
    tileGridSize=(8,8))
enhanced = clahe.apply(gray)
            \end{lstlisting}
        \end{block}

        \column{0.5\textwidth}
        \begin{block}{二值化函数}
            \begin{lstlisting}[language=Python, basicstyle=\ttfamily\tiny]
# 全局阈值
ret, binary = cv2.threshold(
    blur, 127, 255,
    cv2.THRESH_BINARY)

# Otsu自动阈值
ret, binary = cv2.threshold(
    blur, 0, 255,
    cv2.THRESH_BINARY + cv2.THRESH_OTSU)

# 自适应阈值
binary = cv2.adaptiveThreshold(
    blur, 255,
    cv2.ADAPTIVE_THRESH_MEAN_C,
    cv2.THRESH_BINARY, 11, 2)
            \end{lstlisting}
        \end{block}

        \begin{block}{几何矫正函数}
            \begin{lstlisting}[language=Python, basicstyle=\ttfamily\tiny]
# 获取变换矩阵
M = cv2.getPerspectiveTransform(
    src_pts, dst_pts)

# 透视变换
result = cv2.warpPerspective(
    img, M, (w, h))
            \end{lstlisting}
        \end{block}
    \end{columns}
\end{frame}

% -----------------------------------------------------------------------------
% 第3页：课后作业
% -----------------------------------------------------------------------------

\begin{frame}[fragile]{课后作业：完整试卷预处理流水线}
    \begin{alertblock}{作业要求}
        实现一个完整的试卷预处理流水线，处理一张倾斜、模糊、光照不均的试卷照片。
    \end{alertblock}

    \begin{columns}
        \column{0.55\textwidth}
        \textbf{处理流程（必须包含）：}
        \begin{enumerate}
            \item \highlight{去噪}：使用 GaussianBlur 或 medianBlur
            \item \highlight{增强}：使用 CLAHE 对抗光照不均
            \item \highlight{二值化}：使用自适应阈值或 Otsu
            \item \highlight{几何矫正}：使用 warpPerspective 矫正倾斜
        \end{enumerate}

        \vspace{0.3cm}
        \textbf{提交要求：}
        \begin{itemize}
            \item 完整代码（.py 文件）
            \item 处理前后对比图
            \item 简要说明每个步骤的作用
        \end{itemize}

        \column{0.45\textwidth}
        \begin{block}{参考框架}
            \begin{lstlisting}[language=Python, basicstyle=\ttfamily\tiny]
import cv2
import numpy as np

def preprocess_exam(img):
    # 1. 灰度化
    gray = cv2.cvtColor(
        img, cv2.COLOR_BGR2GRAY)

    # 2. 去噪
    blur = cv2.GaussianBlur(
        gray, (5, 5), 0)

    # 3. 增强
    clahe = cv2.createCLAHE(2.0)
    enhanced = clahe.apply(blur)

    # 4. 二值化
    binary = cv2.adaptiveThreshold(
        enhanced, 255,
        cv2.ADAPTIVE_THRESH_MEAN_C,
        cv2.THRESH_BINARY, 11, 2)

    # 5. 几何矫正（需要4个角点）
    # ... warpPerspective

    return binary
            \end{lstlisting}
        \end{block}
    \end{columns}
\end{frame}

% -----------------------------------------------------------------------------
% 第4页：知识点网络与下周预告
% -----------------------------------------------------------------------------

\begin{frame}{知识点网络与下周预告}
    \begin{columns}
        \column{0.55\textwidth}
        \textbf{本周核心知识点：}
        \begin{itemize}
            \item \highlight{去噪}：高斯滤波、中值滤波
            \item \highlight{对比度增强}：CLAHE 自适应直方图均衡
            \item \highlight{二值化}：全局阈值、自适应阈值、Otsu
            \item \highlight{几何矫正}：透视变换 warpPerspective
            \item 完整预处理流水线
        \end{itemize}

        \vspace{0.3cm}
        \textbf{下周预告（Week 4）：试卷版面分析}
        \begin{itemize}
            \item \textbf{边缘检测}：Canny 算子
            \item \textbf{轮廓查找}：findContours
            \item \textbf{轮廓近似}：approxPolyDP
            \item \textbf{定位答题卡}：利用轮廓定位答题区域
        \end{itemize}

        \column{0.45\textwidth}
        \begin{block}{跨周链接}
            \begin{itemize}
                \item Week 1：图像基础
                \item Week 2：AI辅助编程
                \item Week 3：图像预处理 $\checkmark$
                \item Week 4：版面分析 $\triangleright$
                \item Week 5：选择题识别
            \end{itemize}
        \end{block}

        \vspace{0.3cm}
        \begin{alertblock}{重点提示}
            下周我们将进入\textbf{版面分析}，学会如何从整张试卷图像中\highlight{自动定位}出答题卡的位置！
        \end{alertblock}
    \end{columns}
\end{frame}

% -----------------------------------------------------------------------------
% 第5页：总结与问答
% -----------------------------------------------------------------------------

\begin{frame}{总结与问答}
    \begin{center}
        \Huge \textbf{本周总结}
        \vspace{0.8cm}

        \begin{block}{核心收获}
            \begin{itemize}
                \item 学会了图像预处理的完整流水线
                \item 掌握了去噪、增强、二值化、几何矫正四大技术
                \item 能够处理实际中的"脏"试卷图像
            \end{itemize}
        \end{block}

        \vspace{0.8cm}
        \Huge \textbf{Q \& A}
        \vspace{0.5cm}

        \Large 有什么问题吗？
    \end{center}
\end{frame}
